\chapter{Escopo do Sistema de Banco de Dados}
\label{cap:escopo}

Com base na caracterização do problema de zeladoria urbana e em seu alinhamento às metas dos Objetivos de Desenvolvimento Sustentável, este capítulo delimita o escopo de alto nível do sistema de banco de dados relacional proposto. Seguindo a metodologia de projeto de banco de dados preconizada por \citeonline{elmasri2019sistemas} e \citeonline{heuser2009projeto}, apresenta-se aqui a especificação conceitual do minimundo: os subdomínios funcionais, as principais entidades de negócio, as regras de negócio estruturantes e as consultas gerenciais e operacionais que o sistema deve responder. O detalhamento completo dos atributos, chaves identificadoras, cardinalidades no padrão brModelo e o mapeamento para o esquema relacional compõem o escopo específico da Etapa 2 deste projeto.

\section{Visão Geral do Minimundo e Subdomínios Funcionais}
\label{sec:subdominios}

O banco de dados gerencia, de maneira integrada e auditável, todo o ciclo de vida do monitoramento e da manutenção colaborativa da infraestrutura urbana municipal. Para organizar a complexidade do domínio sem perda de coesão, o escopo do sistema estrutura-se em três subdomínios funcionais interdependentes: (i) \textbf{Demandas do Cidadão, Base Geográfica e Avaliação Social}, que abrange a padronização territorial por Código de Endereçamento Postal (CEP), o cadastro dos moradores, a abertura de ocorrências geolocalizadas com arquivos de mídia comprobatória e a avaliação final de satisfação; (ii) \textbf{Estrutura Operacional, Recursos Humanos e Frota Municipal}, que contempla a organização das secretarias e órgãos da prefeitura, a coordenação das equipes de manutenção por categoria e turno, os funcionários operacionais e a frota de veículos; e (iii) \textbf{Execução de Ordens de Serviço, Almoxarifado e Cadeia de Suprimentos}, que compreende o ciclo das Ordens de Serviço (OS), o catálogo e o saldo de materiais de almoxarifado, o consumo de insumos por reparo e o fornecimento por lotes de compra junto a fornecedores.

\section{Principais Entidades de Negócio}
\label{sec:entidades-negocio}

A representação conceitual do sistema fundamenta-se em \textbf{doze entidades de negócio} principais. No primeiro subdomínio, a entidade \textbf{Região de CEP} (\texttt{regiao\_cep}) estabelece a base geográfica padronizada do município (cidade, bairro, logradouro e referência), servindo de âncora espacial tanto para a residência do \textbf{Cidadão} (\texttt{cidadao}) quanto para a localização da \textbf{Ocorrência} (\texttt{ocorrencia}), chamado formal registrado com protocolo único, categoria do defeito, nível de gravidade e status. A entidade \textbf{Anexo de Mídia} (\texttt{anexo\_midia}) armazena os metadados de fotos ou vídeos comprobatórios enviados no relato, enquanto a \textbf{Avaliação de Atendimento} (\texttt{avaliacao\_atendimento}) registra a nota de satisfação (1 a 5), o comentário e a confirmação de resolução após o reparo.

No segundo subdomínio, a entidade \textbf{Secretaria ou Órgão} (\texttt{secretaria\_orgao}) representa cada pasta municipal (Obras, Iluminação, Limpeza), incluindo seu secretário responsável e orçamento mensal. Cada secretaria coordena instâncias de \textbf{Equipe de Manutenção} (\texttt{equipe\_manutencao}), integradas por servidores modelados na entidade \textbf{Funcionário} (\texttt{funcionario}), além de manter os caminhões e utilitários cadastrados em \textbf{Veículo da Frota} (\texttt{veiculo\_frota}). No terceiro subdomínio, a \textbf{Ordem de Serviço} (\texttt{ordem\_servico}) conecta a ocorrência à equipe e ao veículo destacados, registrando cronologia, relatório técnico e custo de mão de obra; \textbf{Material ou Insumo} (\texttt{material\_insumo}) controla o estoque e o custo médio do almoxarifado; e \textbf{Fornecedor} (\texttt{fornecedor}) identifica as empresas que abastecem a prefeitura. A Tabela~\ref{tab:entidades-escopo} sintetiza as doze entidades.

\begin{table}[htb]
\ABNTEXfontereduzida
\caption{Síntese das doze entidades de negócio que compõem o escopo do banco de dados por subdomínio funcional}
\label{tab:entidades-escopo}
\centering
\begin{tabular}{L{3.2cm} L{4.2cm} L{7.0cm}}
\toprule
\textbf{Subdomínio Funcional} & \textbf{Entidade de Negócio} & \textbf{Papel Conceitual no Sistema de Banco de Dados} \\
\midrule
1. Demandas do Cidadão e Base Geográfica & Região de CEP (\texttt{regiao\_cep}) & Padroniza a base territorial (CEP, cidade, bairro, logradouro e referência) para moradias e ocorrências \\
1. Demandas do Cidadão e Base Geográfica & Cidadão (\texttt{cidadao}) & Representa o morador cadastrado que reporta problemas urbanos e avalia a execução dos serviços \\
1. Demandas do Cidadão e Base Geográfica & Ocorrência (\texttt{ocorrencia}) & Registra o chamado de falha na infraestrutura urbana com protocolo, categoria, gravidade e localização \\
1. Demandas do Cidadão e Base Geográfica & Anexo de Mídia (\texttt{anexo\_midia}) & Armazena metadados de fotos e vídeos comprobatórios vinculados existencialmente a uma ocorrência \\
1. Demandas do Cidadão e Base Geográfica & Avaliação de Atendimento (\texttt{avaliacao\_atendimento}) & Registra o parecer final do cidadão (nota de 1 a 5, comentário e confirmação de solução) sobre a OS \\
\midrule
2. Estrutura Operacional e Recursos Municipais & Secretaria ou Órgão (\texttt{secretaria\_orgao}) & Modela os departamentos municipais gestores de orçamento mensal, equipes de campo e frota \\
2. Estrutura Operacional e Recursos Municipais & Equipe de Manutenção (\texttt{equipe\_manutencao}) & Representa os times operacionais especializados por tipo de problema e alocados por turno \\
2. Estrutura Operacional e Recursos Municipais & Funcionário (\texttt{funcionario}) & Cadastra os servidores e técnicos operacionais vinculados às equipes de manutenção \\
2. Estrutura Operacional e Recursos Municipais & Veículo da Frota (\texttt{veiculo\_frota}) & Controla os veículos operacionais das secretarias usados no deslocamento e transporte de materiais \\
\midrule
3. Execução de OS e Cadeia de Suprimentos & Ordem de Serviço (\texttt{ordem\_servico}) & Formaliza a mobilização de equipe e veículo para reparar uma ocorrência, apurando prazos e custos \\
3. Execução de OS e Cadeia de Suprimentos & Material ou Insumo (\texttt{material\_insumo}) & Mantém o catálogo de materiais de almoxarifado, saldo em estoque e custo unitário médio \\
3. Execução de OS e Cadeia de Suprimentos & Fornecedor (\texttt{fornecedor}) & Identifica as empresas contratadas que fornecem lotes de materiais e insumos à administração \\
\bottomrule
\end{tabular}
\legend{Fonte: Elaborado pelos autores em conformidade com a especificação técnica do projeto (2026).}
\end{table}

\section{Regras de Negócio e Restrições de Integridade do Domínio}
\label{sec:regras-negocio}

O banco de dados incorpora regras de negócio rigorosas para assegurar a consistência semântica do domínio: (i) \textbf{localização padronizada por CEP}, pela qual todo cidadão reside em um CEP cadastrado e toda ocorrência referencia obrigatoriamente o CEP do local do problema; (ii) \textbf{autoria e comprovação}, onde cada ocorrência é aberta por um único cidadão e pode possuir múltiplos anexos de mídia com dependência existencial; (iii) \textbf{subordinação operacional}, em que cada secretaria coordena equipes de manutenção (compostas por funcionários) e mantém sua frota de veículos; (iv) \textbf{mobilização em Ordem de Serviço}, permitindo que uma ocorrência gere uma ou mais Ordens de Serviço, cada qual vinculada a uma ocorrência, uma equipe executora e um veículo de deslocamento; (v) \textbf{controle associativo de materiais e compras}, registrando a quantidade utilizada e a data de baixa de cada insumo consumido por Ordem de Serviço, bem como os múltiplos lotes de compra entregues por fornecedores (discriminados por número de lote, data de entrega, quantidade comprada e preço unitário); e (vi) \textbf{avaliação única por atendimento}, assegurando que cada Ordem de Serviço concluída receba no máximo uma avaliação final registrada pelo cidadão autor do chamado.

\section{Consultas e Questões Gerenciais que o Sistema Deve Responder}
\label{sec:consultas-sistema}

A estruturação dessas entidades e regras (que origina, na Etapa 2, quatorze tabelas relacionais a partir das doze entidades fortes e das duas relações associativas de consumo e fornecimento) permite responder a cinco grupos principais de consultas gerenciais e operacionais: (i) \textbf{diagnóstico geográfico de pontos críticos}, identificando bairros e CEPs com maior incidência de falhas por categoria e nível de gravidade; (ii) \textbf{desempenho e tempo de resposta}, apurando o tempo médio entre abertura da ocorrência, emissão da Ordem de Serviço e conclusão do reparo por secretaria e equipe; (iii) \textbf{execução orçamentária e custo real}, somando o custo de mão de obra das Ordens de Serviço ao valor dos materiais consumidos frente ao orçamento mensal de cada secretaria; (iv) \textbf{gestão de estoque e fornecedores}, auditando insumos com estoque crítico, histórico de preços por lote de compra e utilização da frota; e (v) \textbf{auditoria de satisfação cidadã}, calculando a nota média de atendimento por região e a taxa de ocorrências marcadas pelo morador como não resolvidas.
